\documentclass[11pt]{article}
\usepackage[margin=0.75in]{geometry}
\usepackage[T1]{fontenc}
\usepackage[utf8]{inputenc}
\usepackage{lmodern}
\usepackage{microtype}
\usepackage{amsmath}
\usepackage{booktabs}
\usepackage{tabularx}
\usepackage{enumitem}
\usepackage{titlesec}
\usepackage{graphicx}
\usepackage{xcolor}
% Draft-only marker for content that waits on a result; remove every \rework before submitting.
\newcommand{\rework}[1]{\noindent\colorbox{yellow!40}{\parbox{\dimexpr\linewidth-2\fboxsep}{%
  \small\textbf{REWORK:} #1}}\par}
\usepackage[hidelinks]{hyperref}

\setlist{nosep, leftmargin=*}
\titlespacing*{\section}{0pt}{1.0ex}{0.5ex}
\titlespacing*{\subsection}{0pt}{0.8ex}{0.3ex}
\titleformat*{\subsection}{\normalsize\bfseries}
\setlength{\parskip}{0.5ex}
\setlength{\parindent}{0pt}
\newcolumntype{Y}{>{\raggedright\arraybackslash}X}

\begin{document}

\begin{center}
  {\Large\bfseries Lab Report 4: The Integrated Loop under Planted Failures and a Compute Budget}\\[0.5ex]
  Jamie Carpenter (no team) \quad$\cdot$\quad AM/ME/CE 295 \quad$\cdot$\quad 29 November 2026\\
  \url{https://github.com/jamie-carpenter/lattice-gnn}
\end{center}

\noindent Lab Report~3 showed the diagnostician deciding inside a working loop, against a rule
table and against Optuna, on a subset where no search procedure improved on the untuned surrogate
with the positive-definite (PD) output layer. That result left the agent's value untested on the
decisions it was meant for. This report therefore runs the integrated loop on faults whose causes
are known, under a budget that charges each training run its true computational cost, and states
where the system currently fails. The repository README now describes every step from setup to
scoring.

\section*{1 \quad The integrated system}

\subsection*{1.1 \quad What runs end to end}
A single command runs the loop on the GNN surrogate for the homogenized Young's and shear moduli
and the Poisson coupling terms of periodic strut lattices~\cite{lumpe}. In each round, Optuna
trains a fixed number of models over the current search ranges, the evaluation tool writes the
validation errors per target and per lattice family together with the PD check, and code judges
the previous decision against the seed-to-seed spread and checks the budget. The diagnostician
(\texttt{claude-sonnet-5-5}, called without tools) then returns one decision, \emph{retune},
\emph{change structure} or \emph{escalate}, which the validator accepts or rejects before anything
runs. The rule table can replace the agent in the same loop, so the two deciders differ only in
how they decide.

\rework{Update the design figure (\texttt{scripts/make\_design\_figure.py}) to show the
CPU-hour budget and the planted-failure harness, and include it here.}

\subsection*{1.2 \quad Changes since Lab Report 3}
\begin{itemize}
\item \textbf{Compute-aware budget.} Every run now records its CPU time, summed over all threads,
  and its wall time. A study stops when its CPU-hours, its run count, or whichever is set first
  runs out. The agent sees the cost of each past run and the CPU-hours that remain, and its
  instructions state that wider and deeper networks buy fewer runs. CPU time rather than wall
  time is charged because the laptop's standby does not consume it: in Lab Report~3 one run
  logged 39.9\,h of wall time, almost all of it in standby.
\item \textbf{Width cap lifted.} Lab Report~3 capped the width at 256 because a run-count budget
  let the agent choose width 512 at no cost, although such runs took 11--20\,h each in the
  pilot. With compute charged, width 512 is again available to every condition.
\item \textbf{Affordable rounds only.} The validator, and now the loop itself, require that the
  remaining budget pays for a full round at the mean cost per run so far. When it does not, the
  study stops on its budget instead of asking for a decision, so a small remainder is never
  logged as a forced escalation.
\item \textbf{Planted-failure harness.} Training accepts label faults applied to copies of the
  data (Section~2), and a runner executes, repeats and scores the planted cases.
\end{itemize}

\subsection*{1.3 \quad Setting the budget}
On this laptop's eight threads, CPU time accumulates about six times faster than wall time
(0.230 against 0.038\,h in one smoke run). The Optuna study of Lab Report~3 trained for 3.6\,h
over 20 runs, so each study here receives 22 CPU-hours, which keeps the compute comparable
with Lab Report~3. A width-512 run costs ten to fifty times as much as a width-64 run (0.16
against 0.003--0.007 CPU-hours over two epochs), so at this budget a single 100-epoch run at
width 512 consumes a substantial share of a study.

\section*{2 \quad Evaluation on planted failures}

\subsection*{2.1 \quad Design}
The three planned faults of Lab Report~2 map one-to-one onto the agent's actions, so a
three-line script would pass them. Each case below instead injects a fault into the setup of
Section~1.3 such that the one-line rule a script would apply gives the wrong answer, or such that
no answer is correct. The answer key of every case was written into the repository before the
case was run, is checked by code alone (by the decision type, the new search ranges, the menu
option, and whether the decision's text names a family), and never enters the agent's input; a
unit test asserts the last point. Table~\ref{tab:cases} lists the cases.

\begin{table}[!htb]\centering\small
\caption{Planted cases (\texttt{config/planted.yaml}). All start from the PD layer off unless
stated; lr is the learning rate. The first three have a correct remedy, the last three have
none.}\label{tab:cases}
\begin{tabularx}{\linewidth}{@{}lYY@{}}\toprule
Case & Fault & Answer key (checked by code)\\\midrule
\texttt{collapse} & lr range $10^{-2}$--$10^{-1}$ & first decision retunes lr below
  $3\times10^{-3}$; any structural change first is wrong\\
\texttt{units\_low\_lr} & moduli of validation family \texttt{tet\_Z04.8} (20 lattices) scaled
  by 10, as in a units error; lr range $10^{-5}$--$3\times10^{-5}$ & some retune raises the whole
  lr range to $10^{-4}$ or above, and some escalation names \texttt{tet\_Z04.8}\\
\texttt{width\_trap} & widths $\{256, 512\}$ at lr $5\times10^{-4}$--$1.5\times10^{-3}$, where
  width 512 diverged in the pilot; three runs per round & first decision drops width 512 or
  brings lr to $4\times10^{-4}$ or below\\\midrule
\texttt{control} & none; PD layer on, ranges of Lab Report~3's good runs & any structural change
  is unsupported\\
\texttt{hidden\_noise} & labels shuffled among the 126 lattices of the largest training family;
  PD layer on & any structural change is unsupported\\
\texttt{two\_causes} & lr range $10^{-5}$--$5\times10^{-5}$ & enabling the PD layer, raising lr
  and escalating are all defensible; disabling augmentation is unsupported\\\bottomrule
\end{tabularx}
\end{table}

Each case with a known remedy runs three times: the full loop with the agent; ten further agent
calls on that loop's frozen first-round report, each treated as a first decision, which also
measures how consistent the agent's decision is; and the full loop with the rule table. Each case
without a correct answer trains one round and then receives the same ten agent calls and one
rule-table call. The Optuna baseline is deferred to the final report.

\subsection*{2.2 \quad What the probes changed}
Before fixing the cases, I trained single models at the extreme learning rates
(Table~\ref{tab:probe}). They overturned the premise of one planned case. I had expected an
unstable learning rate to produce non-PD predictions that a rule would misattribute to the output
layer. Instead, a learning rate of $3\times10^{-2}$ or $10^{-1}$ collapses the surrogate to
predicting the mean, whose couplings form a PD compliance block, so the non-PD rate is zero.
Non-PD predictions appear only in a partly trained model, as at $10^{-5}$. The planned
\emph{divergence} case therefore became \texttt{collapse}, and \texttt{two\_causes} moved from a
high to a low learning-rate range, where non-PD predictions and under-training coexist.

\begin{table}[!htb]\centering\small
\caption{Learning-rate probes: width 128, four layers, 100 epochs at most, PD layer off,
validation families of the 2,000-lattice subset. Good runs in Lab Report~3 reached a loss of
about 0.22.}\label{tab:probe}
\begin{tabular}{@{}lccccc@{}}\toprule
lr & Loss & Non-PD rate & Median modulus error & Best epoch / epochs run & CPU-hours\\\midrule
$10^{-1}$ & 1.009 & 0 & 583\% & 19 / 39 & 0.53\\
$3\times10^{-2}$ & 1.008 & 0 & 578\% & 4 / 24 & 0.40\\
$5\times10^{-3}$ & 0.873 & 0 & 312\% & 81 / 100 & 0.93\\
$10^{-5}$ & 0.542 & 0.042 & 169\% & 96 / 100 & 0.97\\\bottomrule
\end{tabular}
\end{table}

\section*{3 \quad Results}
\rework{Fill from \texttt{python -m agents.diagnostician.planted score runs/planted-*}: for each
known case, whether the agent loop and the rule-table loop passed each check, and the pass rate
over the ten frozen calls; for each case without an answer, the spread of first decisions and the
number of unsupported ones. Add CPU-hours used, runs, stop reason and best loss per loop, and one
example decision per case from \texttt{decisions.jsonl}. Started 24 November, 14:28.}

\begin{table}[!htb]\centering\small
\caption{Known cases (to be filled from the study).}\label{tab:known}
\begin{tabular}{@{}llll@{}}\toprule
Case & Agent loop & Agent, 10 frozen calls & Rule-table loop\\\midrule
\texttt{collapse} & & & \\
\texttt{units\_low\_lr} & & & \\
\texttt{width\_trap} & & & \\\bottomrule
\end{tabular}
\end{table}

\begin{table}[!htb]\centering\small
\caption{Cases without a correct answer (to be filled from the study).}\label{tab:ambiguous}
\begin{tabular}{@{}lll@{}}\toprule
Case & Agent, first decisions over 10 calls (unsupported) & Rule table\\\midrule
\texttt{control} & & \\
\texttt{hidden\_noise} & & \\
\texttt{two\_causes} & & \\\bottomrule
\end{tabular}
\end{table}

\section*{4 \quad Current failure modes}
The guideline asks where the system breaks, and the most instructive failures so far lie in the
infrastructure and the evaluation rather than in the agent's reasoning.
\begin{itemize}
\item \textbf{A failed call recorded as a decision.} When the Claude Code sign-in expired, every
  agent call returned an authentication error, which the loop parsed as an invalid answer, retried
  once, and then logged as a forced escalation. A study run in this state would have reported
  plausible but fabricated agent behaviour. A failed call now stops the run with an error and is
  never logged as a decision; a unit test covers this case.
\item \textbf{Training is not reproducible on the CPU.} In Lab Report~3 the agent and the rule
  table trained identical settings, yet their losses differed by up to 0.024, which exceeds the
  0.02 threshold that defines an improvement on the subset. Single verdicts near the threshold
  are therefore uninformative, and comparisons need repeats.
\item \textbf{The epoch cap limits accuracy.} Most good runs reached their best epoch at 92--100
  of 100, and the low-lr probe was still improving at epoch 96, but the number of epochs is not
  a setting the agent may change.
\item \textbf{Unchecked evidence.} The validator checks every field that changes the study, but
  not the agent's evidence text, which in the pilot contained arithmetic slips about the
  remaining budget.
\item \textbf{The cost measure itself.} CPU time excludes standby but still depends on the number
  of threads and on contention from other processes, so the budget is comparable across
  conditions run on this laptop, not across machines.
\item \textbf{Scoring by string match.} Whether an escalation names the corrupted family is
  checked by a substring search in the decision's text; an escalation that describes the family
  without naming it would be scored as a miss.
\item \textbf{A wrong premise in the evaluation.} The probes above show that a planted fault can
  fail to produce the symptom it was designed for; every case therefore records the probe that
  supports its premise.
\end{itemize}
\rework{Add the failures the study reveals: wrong or forced decisions, early escalations,
budget spent on width 512, inconsistent frozen calls, and the rule table's failure modes.}

\section*{5 \quad Next steps (to the final report, 9 December)}
\begin{itemize}
\item The Optuna baseline on every planted case, at the same CPU-hour budget.
\item Repeated loops for the cases where the agent and the rule table differ, and one evaluation
  of each condition's best model on the 64 locked test families.
\item The menu options \texttt{family\_balanced\_loss} and \texttt{physics\_residual}.
\item The final presentation and the four-page paper.
\end{itemize}

\section*{Progress log}
\textbf{Done since Lab Report~3} (\texttt{LOG.md}, 24 November): the CPU-hour budget and the
affordable-round rule; planted label faults, six planted cases with code-checked answer keys,
frozen repeats and scoring; the learning-rate probes; the stop on failed agent calls; an updated
README; 94 tests in total.
\textbf{Remaining:} the study results, the Optuna baseline on the planted cases, repeats and the
test-set evaluation. \textbf{Bottleneck:} CPU time; the planted study needs about a day of laptop
time with the machine plugged in.
\rework{Commit count and log dates since Lab Report~3 was submitted.}

\section*{AI usage disclosure}
\textbf{Tools.} Claude Code (\texttt{claude-opus-5-5}) built this project under my direction: it
wrote the code, tests and scripts and drafted this report from \texttt{LOG.md}, the run folders
and the probe results. The diagnostician itself is \texttt{claude-sonnet-5-5}, called through
\texttt{claude -p} with no tools.
\textbf{My role and checks.} I chose the compute-aware budget following the TA's feedback on Lab
Report~3, the budget of 22 CPU-hours, the planted cases including the ones without a correct
answer, and the decision to defer the Optuna baseline. I reviewed each step and each commit.
Every number in this report comes from a logged run in \texttt{runs/} or from \texttt{LOG.md}.

\begin{thebibliography}{9}\footnotesize\setlength{\itemsep}{0pt}
\bibitem{lumpe} T.~S. Lumpe, T.~Stankovi\'c. Exploring the property space of periodic cellular
  structures based on crystal networks. \emph{PNAS} 118(7), 2021. doi:10.1073/pnas.2003504118.
\end{thebibliography}

\end{document}
